\chapter{מבוא: יום אחד בחיים מחוברים}

\begin{epigraph}{}
פעם שאלנו "מה השעה?". היום אנחנו שואלים "כמה סוללה נשארה?". אולי זו השאלה החשובה יותר.
\end{epigraph}

\section{שבע בבוקר, בכל מקום בעולם}

השעון המעורר מצלצל. היד נשלחת אל השידה עוד לפני שהעיניים נפתחות. היא לא מחפשת משקפיים, לא כוס מים ולא את בן הזוג. היא מחפשת את הטלפון.

בתוך שלושים שניות כבר ידוע לכם מה מזג האוויר, מי כתב לכם בלילה, מה קרה בצד השני של כדור הארץ, ואיזה חבר מהתיכון נסע לתאילנד. עוד לא קמתם מהמיטה, ומאות בני אדם כבר נכנסו לחדר השינה שלכם.

הסצנה הזאת חוזרת על עצמה מיליארדי פעמים בכל בוקר. בתל אביב, בלאגוס ובסאו פאולו. אצל סטודנטית בירושלים ואצל נהג מונית בג'קרטה. לפי ההערכה המעודכנת של איגוד התקשורת הבין-לאומי, כ-5.87 מיליארד בני אדם השתמשו באינטרנט ב-2024, כ-\num{72.3\%} מאוכלוסיית העולם \parencite{itu2026data}. מעולם לא היה דבר כל כך משותף לכל כך הרבה בני אדם. לא דת, לא שפה ולא מטבע.

ובכל זאת, כדאי לזכור גם את הצד השני של המספר. אם כ-\num{72\%} מהאנושות מחוברים, פירוש הדבר שכ-\num{28\%} עדיין אינם מחוברים \parencite{itu2026data}. מיליארדי בני אדם חיים בעולם שבו "לבדוק בגוגל" איננו פועל יומיומי. הם לא מתגעגעים לרשת, כי מעולם לא חיו בתוכה. אבל העולם שסביבם, הבנקים, בתי הספר ומשרדי הממשלה, עובר אליה במהירות. הניתוק, שפעם היה מצב טבעי, הופך בהדרגה לסוג חדש של עוני.

וישראל? בתחילת 2025 כ-\num{91.1\%} מהישראלים היו מחוברים לרשת \parencite{kemp2025israel}. במדינה שבה קשה להסכים כמעט על כל דבר, הסמארטפון הוא אולי הקונצנזוס האחרון.

\section{המהפכה השקטה}

מהפכות גדולות נוטות להיות רועשות. למהפכה החקלאית נדרשו אלפי שנים, והיא שינתה את גוף האדם, את המשפחה ואת הנוף. המהפכה התעשייתית מילאה את השמיים בעשן. המהפכה הדיגיטלית, לעומתן, נכנסה לחיינו בשקט כמעט מוחלט. היא לא שרפה ערים ולא הפילה מלכים. היא רק ביקשה מאיתנו ללחוץ על "אני מסכים".

והיא לא נשארה בכיס. לפי הדוח הגלובלי של \en{DataReportal}, משתמש אינטרנט טיפוסי מבלה ברשת חלק ניכר מכל יום \parencite{kemp2024}. חשבון פשוט ממחיש עד כמה.

\begin{factbox}[נתון מפתח]
\textbf{6 שעות ו-40 דקות} ביום ברשת, אצל משתמש אינטרנט טיפוסי \parencite{kemp2024}. אם תכפילו את זה ב-365 ימים, תקבלו כ-2,430 שעות בשנה, קצת יותר ממאה ימים שלמים. במילים אחרות, וזה כבר החשבון שלנו ולא של הדוח: משתמש כזה מבלה בכל שנה יותר משלושה חודשים "בתוך" המסך.
\end{factbox}

אבל המספרים לבדם לא מספרים את הסיפור. השאלה המעניינת איננה כמה זמן אנחנו מבלים מול המסך, אלא מה המסך עושה לנו בזמן הזה. האם הוא כלי, כמו פטיש? או סביבה, כמו עיר? או אולי, וזו האפשרות המטרידה ביותר, הוא סוכן בעל אינטרסים משלו?

דמיינו ציידת-לקטת, הרבה לפני הכתב. כדי לשרוד, היא הייתה צריכה להכיר את הצמחים סביבה ואת עקבות בעלי החיים. הידע שלה, מן הסתם, היה בראש. אנחנו מחזיקים בכיס גישה לכמעט כל הידע האנושי, אבל לא תמיד זוכרים את מספר הטלפון של אמא. זה לא בהכרח רע. זה פשוט שונה. וכדאי להבין את השינוי לפני שמחליטים אם הוא טוב.

יש עוד הבדל, עדין יותר. בלי כתב ובלי דואר, אפשר היה להכיר לעומק רק את מי שחי בסביבה הקרובה. אנחנו מחוברים למאות ולאלפים, אבל רובם מוכרים לנו רק כתמונה קטנה ועיגול ירוק שמסמן "מחובר עכשיו". זו השערה ולא ממצא, אבל קשה להשתחרר ממנה: אם המוח שלנו עוצב בעולם של קהילות קטנות, הוא נדרש היום לעקוב אחרי כפר גלובלי שלם. לא פלא אם לפעמים הוא מתעייף.

\begin{thinkbox}
נסו להיזכר: מתי בפעם האחרונה עברה עליכם שעה שלמה, בהכרה מלאה, בלי לגעת במסך כלשהו? אם קשה לכם להיזכר, שאלו את עצמכם: האם זו הייתה בחירה שלכם?
\end{thinkbox}

\section{מה הספר הזה מנסה לעשות}

הספר הזה לא נכתב כדי להפחיד, וגם לא כדי לחגוג. הוא מנסה לעשות משהו צנוע יותר וקשה יותר: להסתכל על הטכנולוגיה של חיי היומיום מגובה רב, כמו שהיסטוריון מסתכל על אימפריה, ובאותו זמן מקרוב מאוד, כמו שמסתכלים על המסך בשתיים בלילה.

המסע שלפניכם בנוי משבעה פרקים נוספים:

\begin{enumerate}
  \item \textbf{מהטלגרף לסמארטפון}: איך הגענו לכאן, ולמה כל המצאה חדשה בתקשורת נראתה לבני דורה כסוף העולם.
  \item \textbf{אנטומיה של יום דיגיטלי}: לאן באמת הולכות השעות, ולמה בדרום אפריקה מבלים ברשת כמעט פי שניים מאשר בדנמרק.
  \item \textbf{כלכלת הקשב}: מי מרוויח מהגלילה שלכם, ואיך זה עובד.
  \item \textbf{זמן מסך ורווחה}: מה המדע יודע, ובעיקר מה הוא עדיין לא יודע.
  \item \textbf{הבית החכם}: כשהמקרר, הנורה והשואב מתחילים לדבר ביניהם.
  \item \textbf{איזון דיגיטלי}: ארגז כלים מעשי, בלי מוסר ובלי אשמה.
  \item \textbf{סיכום ומבט לעתיד}: מה קורה כשהמכונה מתחילה לדבר בחזרה.
\end{enumerate}

לאורך הדרך תמצאו תיבות צבעוניות: נתוני מפתח, אזהרות, טיפים ושאלות למחשבה. את רוב השאלות לא נפתור. אבל לפעמים שאלה טובה שווה יותר מתשובה מהירה. בעיקר בעידן שבו תשובות מהירות נמצאות במרחק הקלקה אחת.
